\documentclass[11pt,a4paper]{appletmanual}
\usepackage[english]{babel}
\manuallang{en}

\appletname{Labour Income Tax}
\appletsub{A bracketed tax bends the constraint between consumption and
leisure. Marginal rates kink the frontier; average rates break it, and crossing
a threshold leaves the worker strictly worse off.}
\appletdir{labor-tax}

\begin{document}
\manualtitlepage{Applet manual · Introductory Microeconomics · Universidad de los Andes\\
Translated from the GeoGebra applet “Restricción Presupuestaria con Impuesto a la Renta”.\\
MIT licensed. It opens in a browser; there is nothing to install.}

\manualtoc
\thispagestyle{empty}
\clearpage
\setcounter{page}{1}

% ==========================================================================
\section{What this is}

A worker splits a time endowment between working and not working. Devoting a
fraction $L$ to work earns $wL$, pays a tax on that income, and spends what is
left on consumption:
\[
  c(L) \;=\; \frac{m_0 + wL - T(wL)}{p_x},
  \qquad \ell = 1 - L .
\]

The endowment is normalised to $1$, so $L$ is the fraction of time worked and
$\ell$ is leisure. The only thing to choose is $L$.

All the interest is in $T(\cdot)$, the bracketed tax, and in a design decision
that real tax systems make and that classes almost never discuss: whether each
bracket's rate applies \emph{only to the income falling inside that bracket} or
\emph{to all income}. The two schemes sound nearly identical and produce
completely different budget frontiers.

\figher[1.0]{marginal}{Marginal rates. Each bracket's rate taxes only the
income inside it. Net income stays continuous, so the frontier \emph{kinks} but
does not break.}

\begin{amkey}
Under marginal rates the frontier has a \textbf{kink}: the slope changes
abruptly but the line does not break. Under average rates it has a
\textbf{notch}: earning one more unit leaves the worker with less. A kink
attracts optima; a notch pushes them down and pins them just below the
threshold.
\end{amkey}

\subsection{Who it is for}

The consumption-leisure choice, or a public-economics class on tax design. It
assumes indifference curves, a budget constraint and the idea of an optimum.
Without changing anything it also serves to explain why tax data shows people
bunched just below a threshold.

%% Sections shared by every manual, English edition.
%% Pulled in with \input{common-en}. The only thing that differs between
%% applets is \appletfolder, which is also the folder name and the path on the site.

\section{Opening it}

There is nothing to install. The whole application is one \code{index.html}
file: it downloads no libraries, calls no server, and stores nothing outside
your browser.

\subsection{The three ways}

\begin{description}
\item[Double-click the file.] The quickest route. Find
  \code{\appletfolder/index.html} and open it. The browser shows it straight off the
  disk, with an address beginning \code{file://}. It works offline from the
  first moment.
\item[From the web.] If someone has published it, the address ends in
  \code{/\appletfolder/}. Once the page has loaded you can disconnect: it needs the
  network for nothing else.
\item[Served from a folder.] For a class, putting the folder on any static
  server is enough. With Python installed:
  \begin{quote}\code{python3 -m http.server 8000}\end{quote}
  then \code{http://localhost:8000/\appletfolder/} in the browser.
\end{description}

\begin{amnote}
Saving the page with \key{Ctrl}\,+\,\key{S} (or \key{Cmd}\,+\,\key{S} on a
Mac) leaves a copy of your own that still works offline and that can travel on a
memory stick or go out by email.
\end{amnote}

\subsection{Language and theme}

Two pairs of buttons sit at the top right.

\begin{ctrltable}{Control & What it does}
ES / EN & Switches the whole interface between Spanish and English on the spot.
          Nothing is recomputed, so you can switch mid-explanation without
          losing the state you had. \\
Dark / Light & Switches the theme. Dark is the original design; light exists
          for projectors and for printing. \\
\end{ctrltable}

Both choices are remembered in the browser, so next time it opens the way you
left it.

\subsection{Which browser}

Any reasonably recent one: Chrome, Edge, Firefox or Safari, on a computer,
tablet or phone. On a narrow screen the panels stack one above the other
instead of sitting side by side. No account, no permissions, no extensions.


% ==========================================================================
\section{The screen}

\fig[1.0]{interface}{The whole screen: controls on the left, the verdict along
the top, the consumption-leisure plane and the net-income panel below.}

\begin{description}
\item[How it applies.] The two schemes, marginal rates and average rates. This
  is the control that matters; everything else is context.
\item[Brackets.] Three rates $t_0, t_1, t_2$ and two thresholds $s_0, s_1$, all
  on labour income $wL$.
\item[Wage and income.] The wage $w$, non-labour income $m_0$, and the price of
  consumption $p_x$.
\item[Preferences.] Four families and a box to write your own.
\item[Vertical axis.] Whether the plane's axis measures leisure or labour.
\item[Layers.] What gets drawn.
\end{description}

The verdict band at the top names which kind of optimum came out: a tangency, a
kink, the edge of a drop, not working, or working the whole endowment.

\subsection{The two panels}

\begin{description}
\item[Consumption and leisure.] The utility map under a topographic ramp,
  everything unaffordable veiled over, the frontier with its kinks or notches,
  the thresholds, the indifference curve through the optimum, and the optimum.
\item[Net income.] Net against gross income, with the 45° line for reference.
  Preferences play no part here, which is the point: the notch is a property of
  the schedule, not of the worker. The switch changes it to \ui{Utility}, which
  draws $u(c(L), 1-L)$ along the frontier.
\end{description}

\fig[0.62]{notch}{The net-income panel under average rates. At every threshold
net income falls vertically: at the opening values, crossing the first
threshold takes it from $20$ to $12$.}

\subsection{The dark theme}

\fig[0.95]{dark}{Both panels in the dark theme, under average rates.}

% ==========================================================================
\section{The model}

\subsection{The two schemes}

Let $t_0, t_1, t_2$ be the rates and $s_0 \le s_1$ the thresholds, all on
labour income $Y = wL$.

\subsubsection{Marginal rates}

Each rate taxes only the part of income inside its own bracket:
\[
  T(Y) \;=\; t_0\min(Y, s_0)
        \;+\; t_1\max\bigl(0, \min(Y, s_1) - s_0\bigr)
        \;+\; t_2\max(0, Y - s_1).
\]
$T$ is continuous and increasing, and net income $Y - T(Y)$ is continuous and
concave. The budget frontier is a polyline: at each threshold it changes slope
but does not break.

\subsubsection{Average rates}

The bracket you land in sets the rate on \emph{all} income:
\[
  T(Y) \;=\; Y \cdot t_{i(Y)},
  \qquad
  i(Y) = \begin{cases} 0 & Y < s_0\\ 1 & s_0 \le Y < s_1\\ 2 & Y \ge s_1\end{cases}
\]
Here $T$ jumps at every threshold and net income jumps down. At the opening
values ($t_1 = 0.4$), crossing $s_0 = 20$ takes net income from $20$ to $12$:
the worker who earns one unit more keeps eight less.

\subsection{Kink versus notch}

The two names are easy to confuse and the difference is everything.

\begin{ctrltable}{Name & What it is}
Kink & The function is continuous; its derivative jumps. The frontier changes
  slope. An optimum can land exactly there, and usually does. \\
Notch & The function itself is discontinuous. The frontier breaks and there is
  a stretch of points no $L$ reaches. The optimum is pinned just below. \\
\end{ctrltable}

\begin{amnote}
Both generate bunching in the data, but for different reasons: at a kink
because the marginal rate changes; at a notch because crossing it simply makes
the taxpayer worse off.
\end{amnote}

\subsection{How the optimum is found}

Not from a first-order condition. On these frontiers the best point is very
often at a kink, where the slope does not exist and no tangency holds; and
under average rates it sits at the right-hand end of a piece with a strictly
worse point immediately to its right.

What happens instead:

\begin{enumerate}
\item The frontier is split into pieces, one per bracket, with their own
  endpoints.
\item Each piece is scanned densely and its best sample refined by golden
  section.
\item \emph{Every vertex is tested explicitly}, from both sides, each on its own
  bracket's formula.
\item Whichever wins is the optimum, and where it came from is what classifies
  it.
\end{enumerate}

\begin{amwatch}
Each piece carries its own bracket rather than looking it up from the income.
Looking it up makes $c(L)$ single-valued at a threshold, which silently loses
the left-hand limit --- the highest point of the frontier --- and left it
falling outside the computed view.
\end{amwatch}

% ==========================================================================
\section{The controls}

\subsection{Brackets, wage and income}

\begin{ctrltable}{Slider & What it is}
$t_0$, $t_1$, $t_2$ & The three rates, 0 to 0.95. \\
$s_0$, $s_1$ & The two thresholds, on labour income. $s_1$ is held at or above
  $s_0$: the other way round would invert the schedule. \\
$w$ & The wage. It sets where the thresholds land on the labour axis,
  $L = s/w$. \\
$m_0$ & Non-labour income. Shifts the whole frontier upward. \\
$p_x$ & Price of consumption. Bounded away from zero. \\
\end{ctrltable}

\subsection{Preferences}

\begin{ctrltable}{Family & Form}
Cobb--Douglas & $x^{\alpha} y^{1-\alpha}$ \\
CES & $(\alpha x^{\rho} + (1-\alpha) y^{\rho})^{1/\rho}$ \\
Quasilinear & $x + 4\alpha\ln y$ \\
Custom & Whatever you type \\
\end{ctrltable}

In the expression $x$ is consumption and $y$ is leisure. With $\rho$ strongly
negative the CES approaches perfect complements and corner optima appear.

\subsection{Vertical axis and layers}

The \ui{Vertical axis} switch moves between \ui{Leisure} and \ui{Labour}.
Utility is always read on leisure, so flipping the axis genuinely re-reads
preferences rather than mirroring a picture.

\fig[0.62]{labour-axis}{The plane with the vertical axis set to labour.}

\begin{ctrltable}{Layer & What it draws}
Utility map & The field $u$ under a colour ramp. \\
Background contours & Level curves of utility, independent of the optimum. \\
Feasible set & Veils what cannot be afforded. \\
Curve through the optimum & The indifference curve passing through it. \\
Optimum & The chosen point. \\
Thresholds & The verticals at $L = s_0/w$ and $L = s_1/w$. \\
Grid & The background rule. \\
\end{ctrltable}

% ==========================================================================
\section{The verdict}

\begin{ctrltable}{Verdict & What it means}
Interior optimum & Tangency: the marginal rate of substitution equals the
  bracket's net wage. The worker is inside a bracket. \\
Optimum at the kink & It lands exactly on a threshold. To the left another hour
  still repays; to the right it does not. \\
Optimum pinned below the drop & Under average rates. The next hour would cross
  the threshold and levy the higher rate on all income. \\
Not working is best & Corner at $L = 0$: living on non-labour income. \\
Working the whole endowment & Corner at $L = 1$. \\
\end{ctrltable}

\fig[0.85]{kink-verdict}{The verdict at the opening values under marginal
rates: the optimum lands on the first threshold.}

% ==========================================================================
\section{Classroom activities}

\begin{activity}{The same tax, two frontiers}
\amset Opening values. Switch between \ui{Marginal rates} and \ui{Average
rates} while watching the left panel.
\amexp Under marginal rates the frontier bends at the thresholds. Under average
rates it \emph{breaks}: two stretches of line that do not touch, and a whole
band of consumption levels no amount of work reaches.
\par\smallskip
The rates, thresholds, wage and preferences are identical in both. All that
changed was the phrase ``on the income in the bracket'' to ``on all income''.
\end{activity}

\begin{activity}{Bunching just below}
\amset Opening values, \ui{Marginal rates}.
\amexp The \ui{Labour L*} readout gives $0.400$, which at $w = 50$ is exactly
gross income $20$: the first threshold. The verdict says ``Optimum at the
kink''.
\par\smallskip
Move $t_1$ up and down. The optimum stays pinned to the threshold across a wide
range of rates, because at a kink there is not one slope to match but a whole
interval of them. That is what produces bunching in the data, and it is not
evasion: it is the optimum.
\end{activity}

\begin{activity}{A notch that costs welfare}
\amset \ui{Quasilinear} preferences. Compare \ui{Marginal rates} with
\ui{Average rates}.
\amexp Under marginal: $L^\ast = 0.800$, consumption $38.00$, gross income
$40.0$, tax $12.0$, utility $34.781$. The average rate is $30\%$ and the
marginal rate $80\%$.
\par\smallskip
Under average rates the same worker retreats to $L^\ast = 0.400$, consumption
$30.00$, tax $0.0$, utility $28.978$. It collects \emph{less} tax and the
worker is \emph{worse off}. Ask who that design benefits. Nobody: it is pure
loss, and it is why almost no tax system uses average rates.
\end{activity}

\begin{activity}{Measuring the drop}
\amset \ui{Average rates}, right panel on \ui{Tax}.
\amexp The \ui{Largest drop} readout gives $-12.00$. Find it in the picture:
it is at the second threshold, $s_1 = 30$, where the rate goes from $0.4$ to
$0.8$ on all income. Just below, net income is $30 - 12 = 18$; just above,
$30 - 24 = 6$.
\par\smallskip
The first threshold drops too, from $20$ to $12$, but by less. Lower $t_2$
until the first becomes the largest and watch the readout change.
\end{activity}

\begin{activity}{Utility breaks into pieces}
\amset \ui{Average rates}, right panel on \ui{Utility}.
\amexp The curve $u(c(L), 1-L)$ appears split into disconnected arcs, one per
bracket. Its highest point is at the right-hand end of one of them, not at an
interior maximum.
\par\smallskip
It is the same information as the left panel, one dimension down. Ask where a
student who differentiated and set to zero would look, and why they would not
find it.
\end{activity}

\begin{activity}{Not working}
\amset \ui{CES} preferences, $\rho = -6$.
\amexp The optimum goes to $L^\ast = 0.000$: consumption $10.00$, which is
exactly $m_0/p_x$. The verdict says ``Not working is best''.
\par\smallskip
With $\rho$ strongly negative the CES approaches perfect complements and
leisure becomes nearly indispensable. Raise $w$ and see how far it has to go
before the first hour is worth it.
\end{activity}

\begin{activity}{The axis is not the preference}
\amset Opening values. Move \ui{Vertical axis} from \ui{Leisure} to
\ui{Labour}.
\amexp The utility map is not the mirror image of the previous one. Utility is
always read on leisure, so flipping the axis genuinely changes what is being
drawn, not just how it is being looked at.
\par\smallskip
A useful check against the habit of reading the two axes as interchangeable.
\end{activity}

\begin{activity}{Designing a tax without notches}
\amset \ui{Average rates}. The goal: a frontier that does not break.
\amexp There is only one way: make all three rates equal. Any difference
between $t_0$, $t_1$ and $t_2$ produces a notch at the corresponding threshold.
\par\smallskip
Under marginal rates, by contrast, any increasing combination gives a
continuous frontier. There is the whole design argument, and it fits in a
sentence.
\end{activity}

% ==========================================================================
\section{Expression syntax}

In the utility box, $x$ is consumption and $y$ is leisure.

\begin{ctrltable}{Element & What is allowed}
Operators & \code{+ - * / \^{} ( )} \\
Functions & \code{sqrt ln log log10 exp abs sin cos tan min max pow} \\
Constants & \code{e}, \code{pi} \\
Implicit multiplication & \code{2xy} is \code{2*x*y} \\
Associativity & \code{\^{}} binds to the right and tighter than unary minus:
  \code{-x\^{}2} is $-(x^2)$ \\
\end{ctrltable}

Partial derivatives are computed symbolically. With \code{min}, \code{max} or
\code{abs} --- precisely the kinked cases --- it falls back to central
differences.

% ==========================================================================
\section{Changes from the GeoGebra original}

\begin{description}
\item[One tax base.] \code{TAX\_LABOR} taxed $wL$, but the average-rate branch
  \code{l1} taxed $wL + I_0$. The two schemes were therefore not comparable at
  exactly the moment you want to compare them. Both tax labour income here.
\item[Discontinuities are modelled.] The original wrote the schedule as nested
  \code{If[...]}, which GeoGebra draws with a segment bridging the gap. Each
  bracket is a separate piece here, so a notch stays a notch.
\item[Each piece carries its own bracket.] Looking the bracket up from the
  income makes $c(L)$ single-valued at a threshold, which silently loses the
  left-hand limit --- the highest point of the frontier. It was falling outside
  the computed view and being clipped.
\item[Preferences and an optimum.] The original drew the budget set only.
  Without an optimum the kink and the notch are shapes rather than predictions,
  so a utility function, the indifference curve through the optimum and the
  bunching verdict are added.
\item[$p_X$ and $w$ started at 0], dividing by zero throughout. Both are
  bounded away from it.
\item[$s_1 < s_0$] would invert the schedule; the second threshold is held at
  or above the first.
\item[Dead objects.] \code{s\_2} had a slider and a label but never entered
  \code{TAX\_LABOR}; \code{a}, \code{b}, \code{f}, \code{ngreso}, \code{E},
  \code{E\_trace}, \code{Pendiente}, \code{B\_OCIO}, \code{B\_TRABAJO},
  \code{h\_TRABAJO} and the \code{Q}/\code{P} point set were leftovers from the
  applet this was copied from. Not carried over.
\end{description}

% ==========================================================================
\section{Troubleshooting}

\begin{ctrltable}{Symptom & What is happening}
The plane is blank & The utility expression does not parse. The red message
  under the box says whether it is syntax or a variable. \\
The thresholds are not visible & They are off the axis: with a small $w$,
  $s/w$ can exceed 1. Raise $w$ or lower the thresholds. \\
Both schemes give the same answer & All three rates are equal. With equal rates
  there is neither kink nor notch and the two schemes coincide. \\
Largest drop reads ``---'' & You are on marginal rates, where by construction
  there are no drops. \\
The optimum does not move when I change a rate & It is pinned at a kink. That
  is the correct behaviour and the subject of activity 2. \\
\end{ctrltable}

%% Sections shared by every manual, English edition.
%% Pulled in with \input{common-en}. The only thing that differs between
%% applets is \appletfolder, which is also the folder name and the path on the site.

\section{Publishing it for a class}

Any static host will do: GitHub Pages, Netlify, a university web directory,
even a shared folder served over HTTP. The only thing to upload is the
\code{index.html} file, inside a folder named however you want the address to
read.

The repository ships a GitHub Actions workflow
(\code{.github/workflows/pages.yml}) that publishes every applet at once on each
push to the default branch. To switch it on, in the repository:
\emph{Settings\,$\rightarrow$\,Pages\,$\rightarrow$\,Source: GitHub Actions}.
That is the one step the workflow cannot do for itself, because creating a Pages
site needs administration rights that a \code{GITHUB\_TOKEN} is never granted.

\begin{amnote}
To hand the applet to students without reliable internet, send them the
\code{index.html} file directly. It weighs less than a photograph and opens with
a double-click.
\end{amnote}

\section{Licence and provenance}

The application code is MIT. Use it, change it and pass it on, in class or out
of it, with attribution.

This applet is a standalone rewrite of an original GeoGebra applet. The rewrite
is not a literal translation: where the original had a construction error, a
division by zero reachable with its own sliders, or a dead object inherited from
whatever applet it was copied from, this version fixes it and says so. The
section \emph{Changes from the GeoGebra original} lists those changes one by
one.

Everything is hand-written and dependency-free: the expression parser, the
symbolic differentiation, the level-curve tracing and the drawing. In
particular there is no \code{eval}, so what a student types into a text box
cannot become code, and a strict content security policy does not break the
page.


\end{document}
